\documentclass[ECP,preprint]{ejpecp}

\SHORTTITLE{Variance and local log-concavity of Bernoulli sums}

\TITLE{Variance and local log-concavity of Poisson--binomial laws%
  \thanks{Generative AI systems assisted with this work; see the
    acknowledgements at the end.}}

\AUTHORS{%
  Brett~Reynolds\footnote{Humber College, Toronto, Canada.
    \EMAIL{brett.reynolds@humber.ca}}%
  \orcid{0000-0003-2407-9448}}

\KEYWORDS{Poisson--binomial law; Bernoulli sum; log-concavity; Tur\'an
  inequality; ultra-log-concavity; modal index; random unitary matrix;
  computer-assisted proof}

\AMSSUBJ{60E15}
\AMSSUBJSECONDARY{60C05; 05A20}

\ABSTRACT{Let $W$ be a finite sum of independent Bernoulli variables with
  probability mass function $f$ and variance $V\geq1$, and let $D$ be the
  first index at which $f$ decreases. The normalized Tur\'an deficit
  $\delta_k=1-f_{k-1}f_{k+1}/f_k^2$ is a bounded transform of the discrete
  curvature of $\log f$ at $k$, and for a normal density of variance $V$ that
  curvature is $1/V$. We prove
  $V\delta_D\geq1/4$, and an explicit family of binomial laws shows that no
  constant above $1/3$ is possible, so the best constant lies between $1/4$
  and $1/3$. Combined with cubic inequalities of Hillion and Johnson, this
  gives $\delta_k\geq1/(4V+|k-D|)$ at every $k$ in the support, and at the
  rightmost mode $c$ we have $1/(4V+1)\leq\delta_c<2/V$. Ultra-log-concavity
  alone gives no variance-scaled lower bound on $\delta_D$. For the number of eigenvalues of an $N\times N$
  Haar unitary matrix in a half circle, the bound at $D$ is of order
  $1/\log N$, against order $1/N$ from ultra-log-concavity. The proof combines
  the cubic inequalities with a maximal-mass bound of Bobkov, Marsiglietti, and
  Melbourne to reduce the main inequality to a one-variable inequality, which
  is proved by Bernstein expansions, with a computer-assisted step in exact
  rational arithmetic.}

\hypersetup{
  pdftitle={Variance and local log-concavity of Poisson-binomial laws},
  pdfauthor={Brett Reynolds},
}

\newcommand{\Pp}{\mathbb{P}}
\newcommand{\Var}{\operatorname{Var}}
\newcommand{\Ber}{\operatorname{Bernoulli}}

\begin{document}

\section{Introduction and main results}\label{sec:introduction}

Let
\[
  W=\sum_{i=1}^n B_i,
  \qquad B_i\sim\Ber(p_i),
\]
where the summands are independent. The law of $W$ is Poisson--binomial;
see \cite{tangtang2023} for a recent survey.

Summands with $p_i=0$ can be deleted, and summands with $p_i=1$ only
translate the law, which changes nothing below except by a shift of index
(Section~\ref{sec:reduction}). We therefore assume throughout that every
success probability satisfies $0<p_i<1$. Put
\[
  f_k=\Pp(W=k),\qquad V=\Var(W)=\sum_{i=1}^n p_i(1-p_i),
\]
and call $(f_k)_{k=0}^n$ the probability mass function (pmf) of $W$. Extend
the pmf by $f_{-1}=f_{n+1}=0$. The probability-generating polynomial
$\sum_{k=0}^n f_kx^k$ has only negative real zeros, so, by Newton's
inequalities \cite[equation~(1)]{pitman1997}, the pmf is log-concave: $f_k^2\geq
f_{k-1}f_{k+1}$ for every $k$, an inequality of Tur\'an type. For
$0\leq k\leq n$, we measure its slack relative to $f_k^2$ by the normalized
Tur\'an deficit
\begin{equation}\label{eq:deficit}
  \delta_k
  :=1-\frac{f_{k-1}f_{k+1}}{f_k^2}.
\end{equation}
We have $0\leq\delta_k\leq1$ and $\delta_0=\delta_n=1$.

At an index with
$f_{k-1}f_{k+1}>0$,
\[
  \delta_k=1-e^{-\mathcal C_k},
  \qquad
  \mathcal C_k:=2\log f_k-\log f_{k-1}-\log f_{k+1},
\]
so $\delta_k$ is a bounded transform of the discrete curvature of $\log f$
at $k$. For a normal density of variance $V$, the corresponding curvature is
$1/V$ at every point, which suggests $1/V$ as the natural scale for
$\delta_k$ near the mode. Since $\delta_k$ depends on three adjacent masses
and $V$ on the whole law, we ask whether every Bernoulli sum satisfies
$\delta_k\geq\kappa/V$ near the mode for a universal $\kappa>0$, with no
dependence on the number of summands or on the success probabilities.

The number of summands is the wrong scale for this question. Summands with
$p_i$ near $0$ or $1$ increase $n$ but add little to $V$. The lower bounds on
$\delta_k$ that we know of are not variance-scaled. They depend on $n$, such
as the ultra-log-concavity bound recorded by Hillion and Johnson
\cite{hillionjohnson2016}; on the success probabilities, such as a bound of
Johnson \cite{johnson2017discrete}; or on $k$ alone, such as the bound
$\delta_k>1/(k+1)$ that follows from the strict decrease of
$(k+1)f_{k+1}/f_k$ proved by D\"umbgen and Wellner
\cite[Proposition~2]{duembgenwellner2020}. Pitman's two-sided bounds on
adjacent mass ratios through exponentially tilted laws
\cite[equation~(20)]{pitman1997} also give a positive lower bound on
$\delta_k$. All four degenerate near the mode when
summands with $p_i$ near $0$ and summands with $p_i$ near $1$ accumulate at
fixed variance (see the discussion after Proposition~\ref{prop:ulc}).
Darroch's theorem that every mode lies at distance less than one from the
mean \cite[Theorem~4]{darroch1964} and the maximal-mass bounds of
\cite{bailloncominettivaisman2016,bobkovmarsigliettimelbourne2022} locate the
mode or bound the largest mass, but by themselves give no lower bound on
$\delta_k$.

We state the bound at the first index where $f$ decreases,
\begin{equation}\label{eq:first-descent}
  D:=\min\{1\leq k\leq n:f_k<f_{k-1}\}.
\end{equation}
When $V\geq1$, this set is nonempty and $2\leq D\leq n$
(Section~\ref{sec:reduction}), and, by log-concavity, $c:=D-1$ is the
rightmost modal index.

\begin{theorem}\label{thm:main}
  Under the standing notation above, suppose $V\geq1$ and let $D$ be defined
  by \eqref{eq:first-descent}. Then
  \begin{equation}\label{eq:main-bound}
    V\delta_D
    =V\left(1-\frac{f_{D-1}f_{D+1}}{f_D^2}\right)
    \geq\frac14.
  \end{equation}
\end{theorem}

Without some lower bound on $V$, no positive constant is possible. If $W\sim\Ber(p)$ with $0<p<1/2$, then $D=1$, $\delta_D=1$, and
$V\delta_D=p(1-p)\to0$ as $p\downarrow0$.

Let $\kappa_\star$ denote the infimum of $V\delta_D$ over all finite
Poisson--binomial laws with $V\geq1$, where $D$ and $\delta_D$ are defined
from each law's pmf as above. The theorem and the next proposition give
\begin{equation}\label{eq:constant-bracket}
  \frac14\leq\kappa_\star\leq\frac13.
\end{equation}

\begin{proposition}\label{prop:one-third}
  For every $n\geq5$, there exists a binomial law of variance one whose pmf has
  first-descent index $D=2$ and for which
  \[
    \delta_D=\frac{n+1}{3(n-1)}.
  \]
\end{proposition}

\begin{proof}
  Set
  \[
    p_n=\frac{1-\sqrt{1-4/n}}2,
    \qquad W_n\sim\operatorname{Bin}(n,p_n).
  \]
  Then $np_n(1-p_n)=1$. Write $f=(f_k)$ for the pmf of $W_n$ and put
  $q_k=f_k/f_{k-1}$. Then
  \[
    q_k=\frac{n-k+1}{k}\frac{p_n}{1-p_n}.
  \]
  Since $p_n(1-p_n)=1/n$, we have $p_n>1/n$, so $q_1=np_n/(1-p_n)>1$.
  Both $p_n$, the smaller root of $x(1-x)=1/n$, and $2/(n+1)$ lie in
  $(0,1/2)$, where $x\mapsto x(1-x)$ is increasing. For $n\geq5$ we have
  $n^2-4n-1>0$, which is equivalent to
  \[
    \frac{2}{n+1}\left(1-\frac{2}{n+1}\right)
    =\frac{2(n-1)}{(n+1)^2}>\frac1n,
  \]
  so $p_n<2/(n+1)$. This is equivalent to $q_2<1$, so $D=2$. Finally,
  \[
    \delta_2=1-\frac{q_3}{q_2}
    =1-\frac{2(n-2)}{3(n-1)}
    =\frac{n+1}{3(n-1)}\longrightarrow\frac13.
  \]
\end{proof}

We do not know which, if either, endpoint of \eqref{eq:constant-bracket}
equals $\kappa_\star$. We are not aware of an earlier proof that
$\kappa_\star>0$.

Cubic inequalities of Hillion and Johnson \cite{hillionjohnson2016} for
Bernoulli sums imply that $1/\delta_k$ changes by at most one per lattice step
(Section~\ref{sec:propagation}), so the bound at $D$ gives a bound at every
point of the support.

\begin{corollary}\label{cor:support}
  Under the hypotheses of Theorem~\ref{thm:main}, for $0\leq k\leq n$,
  \begin{equation}\label{eq:support-bound}
    \delta_k\geq\frac1{4V+|k-D|}\geq\frac1{4V+|k-\mathbb EW|+2},
  \end{equation}
  and
  \begin{equation}\label{eq:adjacent-ratio-drop}
    V\left(1-\frac{f_{D+1}}{f_D}\right)\geq\frac14.
  \end{equation}
\end{corollary}

\begin{proof}
  By the reciprocal bound \eqref{eq:reciprocal} of
  Section~\ref{sec:propagation} and Theorem~\ref{thm:main},
  $1/\delta_k\leq1/\delta_D+|k-D|\leq4V+|k-D|$. Darroch's theorem gives
  $|c-\mathbb EW|<1$, so $|D-\mathbb EW|<2$, which gives the second
  inequality. For \eqref{eq:adjacent-ratio-drop}, write
  $\eta_-=f_D/f_{D-1}<1$ and $\eta_+=f_{D+1}/f_D$; then
  $1-\eta_+\geq1-\eta_+/\eta_-=\delta_D$.
\end{proof}

On a window $|k-\mathbb EW|\leq\alpha\sqrt V$, the bound
\eqref{eq:support-bound} stays on the $1/V$ scale for fixed $\alpha$.
Section~\ref{sec:reduction} extends Theorem~\ref{thm:main} and
Corollary~\ref{cor:support} to differences of independent Bernoulli sums.

At the rightmost mode, $\delta_c$ is within constant factors of $1/V$; since
$V\geq1$, the next proposition gives $1/5\leq V\delta_c<2$.

\begin{proposition}\label{prop:mode}
  If $V\geq1$, then
  \[
    \frac1{4V+1}\leq\delta_c<\frac2V.
  \]
\end{proposition}

\begin{proof}
  The lower bound is \eqref{eq:support-bound} with $k=c$. For the upper
  bound, the idea is that exponential tilting moves the mean at a rate equal
  to the tilted variance: the tilts at which $c-1,c$ and then $c,c+1$ are the
  modes differ by $-\log(1-\delta_c)$, and by Darroch's theorem the mean moves
  by less than $2$ between them. For
  $t\in\mathbb R$, let $f^{(t)}_k=e^{tk}f_k/\sum_je^{tj}f_j$, the pmf of the Bernoulli sum
  with parameters $p_ie^t/(1-p_i+p_ie^t)$, and write $\mu(t)$ and $V(t)$ for
  its mean and variance, so $\mu'(t)=V(t)$. The tilted variance of a summand
  is $p_i(1-p_i)e^t/(1-p_i+p_ie^t)^2\geq e^{-|t|}p_i(1-p_i)$, so
  $V(t)\geq e^{-|t|}V$. Since $2\leq D\leq n$, the ratios $q_c=f_c/f_{c-1}\geq1$
  and $q_{c+1}=f_{c+1}/f_c<1$ are defined; put $t_-=-\log q_c\leq0$ and
  $t_+=-\log q_{c+1}>0$. The pmf is strictly log-concave (by \eqref{eq:ulc}
  below), so $f^{(t_-)}$ has exactly the two modes $c-1$ and $c$, and
  $f^{(t_+)}$ exactly the two modes $c$ and $c+1$. By Darroch's theorem
  \cite[Theorem~4]{darroch1964}, applied to these tilted laws (which are again Poisson--binomial),
  $\mu(t_-)\in(c-1,c)$ and $\mu(t_+)\in(c,c+1)$. Hence
  \[
    2>\mu(t_+)-\mu(t_-)=\int_{t_-}^{t_+}V(t)\,dt\geq V(2-e^{t_-}-e^{-t_+})
    =V\Bigl(2-\frac1{q_c}-q_{c+1}\Bigr)\geq V\Bigl(1-\frac{q_{c+1}}{q_c}\Bigr),
  \]
  where the last step uses $(1-1/q_c)(1-q_{c+1})\geq0$, and
  $1-q_{c+1}/q_c=\delta_c$.
\end{proof}

A pmf $(f_k)_{k=0}^n$ is ultra-log-concave of order $n$, written
$\operatorname{ULC}(n)$, if $f_k/\binom nk$ is log-concave in $k$.
Poisson--binomial pmfs are $\operatorname{ULC}(n)$, and this gives
\begin{equation}\label{eq:ulc}
  \delta_k\geq
  \frac{n+1}{(k+1)(n-k+1)};
\end{equation}
see \cite[Definition~3.11 and equation~(41)]{hillionjohnson2016} and, for a
recent treatment via relative log-concavity,
\cite[the discussion following Definition~1.3]{marsigliettimelbourne2026}.
No variance-scaled lower bound on $\delta_D$ holds for $\operatorname{ULC}$
laws in general.

\begin{proposition}\label{prop:ulc}
  For $m\geq1$, let $g=g^{(m)}$ be the pmf on $\{0,\ldots,2m\}$ with
  $g_{m+j}\propto\binom{2m}{m+j}2^{-|j|}$. Then $g$ is
  $\operatorname{ULC}(2m)$ and strictly log-concave, its first-descent index
  is $m+1$, its deficit there is $(2m+1)/(m(m+2))$, and its variance $V_m$
  tends to $4$. Hence $V_m\delta_{m+1}\to0$.
\end{proposition}

\begin{proof}
  The ratio $g_{m+j}/\binom{2m}{m+j}\propto2^{-|j|}$ is log-concave, and the
  product of the strictly log-concave binomial row with it is strictly
  log-concave. The pmf is symmetric about $m$, so $m$ is its unique mode and
  $m+1$ its first descent. There the geometric factors cancel:
  $\delta_{m+1}=1-\binom{2m}{m}\binom{2m}{m+2}\binom{2m}{m+1}^{-2}
  =(2m+1)/(m(m+2))$. Finally, $\binom{2m}{m+j}/\binom{2m}{m}$ increases to
  $1$ as $m\to\infty$, for each fixed $j$. The mean is $m$ by symmetry, so by dominated
  convergence with the envelope $(1+j^2)2^{-|j|}$ the variance converges to
  that of the law on $\mathbb Z$ proportional to $2^{-|j|}$, which is $4$.
\end{proof}

For $g^{(m)}$, $1/\delta_m\leq4/3$ while $1/\delta_{m+1}\to\infty$, so
$1/\delta$ changes by an unbounded amount in one step; by
\eqref{eq:reciprocal}, this cannot happen for a Bernoulli sum.

The bound \eqref{eq:ulc} also degenerates for Bernoulli sums themselves, and
so do the other lower bounds on $\delta_k$ cited above. For $m\geq2$, take
$m$ success probabilities equal to $\varepsilon_m$ and $m$ equal to
$1-\varepsilon_m$, where $\varepsilon_m\leq1/2$ solves
$2m\varepsilon_m(1-\varepsilon_m)=1$. The law has $V=1$, and its pmf is
symmetric about $m$ and strictly log-concave, so $D=m+1$. Johnson's bound
\cite[Definition~1.2, Lemma~5.1, and Example~5.2(2)]{johnson2017discrete}
reads $\delta_k\geq(f_{k+1}/f_k)\bigl(\sum_ip_i/(1-p_i)\bigr)^{-1}$, and here
the odds sum is at least $m(1-\varepsilon_m)/\varepsilon_m\sim2m^2$.
Combining Pitman's \cite[equation~(20)]{pitman1997} at $k-1$ and $k$ gives
$\delta_k\geq1-\theta(k-1/(n-k+2))/\theta(k+1/(k+2))$, where $\theta(x)=e^t$
for the tilt $t$ at which the tilted law of the proof of
Proposition~\ref{prop:mode} has mean $x$. In this family the tilted mean is
$m+\sinh t/(1+(\cosh t-1)/m)$; at $k=D$ the two arguments differ by
$2/m+O(m^{-2})$ near the point $t=\operatorname{arsinh}1$, where the tilted
mean is $m+1+O(1/m)$ and has slope $\sqrt2+O(1/m)$. Since $\log\theta$ is the
inverse of the tilted mean, the two values of $\log\theta$ differ by
$\sqrt2/m+O(m^{-2})$, and so does the bound.
Finally, $W-m$ converges in law to the difference of two independent
Poisson$(1/2)$ variables, whose pmf is $e^{-1}I_k(1)$ with $I_k$ the modified
Bessel function. Table~\ref{tab:balanced} collects the values at $D$.

\begin{table}[ht]
  \centering
  \caption{Lower bounds on $\delta_D$, and $\delta_D$ itself, in the balanced
  family with $V=1$ and $D=m+1$, as $m\to\infty$.}
  \label{tab:balanced}
  \begin{tabular}{lll}
    \hline
    Bound & Source & Value at $D$\\
    \hline
    Ultra-log-concavity \eqref{eq:ulc} & \cite{hillionjohnson2016} &
      $(2m+1)/(m(m+2))\sim2/m$\\
    $\delta_k>1/(k+1)$ & \cite{duembgenwellner2020} & $1/(m+2)\sim1/m$\\
    Odds-sum bound & \cite{johnson2017discrete} & $O(m^{-2})$\\
    Combined tilt bounds & \cite{pitman1997} & $\sim\sqrt2/m$\\
    Theorem~\ref{thm:main} & & $1/(4V)=1/4$\\
    \hline
    $\delta_D$ itself & & $\to1-I_0(1)I_2(1)/I_1(1)^2\approx0.46$\\
    \hline
  \end{tabular}
\end{table}

For every Bernoulli sum with $V\geq1$, \eqref{eq:adjacent-ratio-drop} gives
$f_{D+1}/f_D\leq1-1/(4V)$. Pitman's ratio bound
\cite[equation~(21)]{pitman1997}, with the tilted
variance $V(t)\leq e^tV$ for $t\geq0$ (notation as in the proof of
Proposition~\ref{prop:mode}), gives $f_{k+1}/f_k<V/(V+k-\mathbb EW)$ for
integers $\mathbb EW<k<n$: there $f_{k+1}/f_k<e^{-t_k}$, where $\mu(t_k)=k$,
and $k-\mathbb EW=\int_0^{t_k}V(t)\,dt\leq V(e^{t_k}-1)$. At $k=D$ this bound
tends to one when $D-\mathbb EW\downarrow0$, as for $\operatorname{Bin}(n,p)$
with $np=\ell+1-\varepsilon$ for an integer $\ell\geq1$ and $p<\varepsilon$,
where $D=\ell+1$.

A polynomial with nonnegative coefficients and only real zeros is, after
normalization, the generating function of a Bernoulli sum
\cite[Proposition~1]{pitman1997}, so Theorem~\ref{thm:main} and its
consequences apply to the coefficient sequences of such polynomials, for
example rows of the Stirling and Eulerian arrays, whenever the associated
law has variance at least one. Determinantal point processes supply
Bernoulli sums in which $V$ and $n$ differ in order.

\begin{example}\label{ex:cue}
  Let $X_N$ count the eigenvalues of an $N\times N$ Haar unitary matrix in
  the open upper half of the unit circle, with $N$ even. Meckes and Meckes
  \cite[Proposition~3(1)]{meckesmeckes2017} show that $X_N$ has the law of a
  sum of $N$ independent Bernoulli variables. By the theorem of Hough, Krishnapur, Peres,
  and Vir\'ag on which this rests \cite[Theorem~7]{houghkrishnapurperesvirag2006},
  the parameters are the eigenvalues of the kernel of the eigenvalue process,
  a projection of rank $N$, compressed to the half circle. These are the
  nonzero eigenvalues of the Gram matrix $\Gamma$ of the functions
  $e^{ij\phi}$, $0\leq j<N$, on the half circle, with entries
  $\Gamma_{jk}=(2\pi)^{-1}\int_0^\pi e^{i(j-k)\phi}\,d\phi$. Both
  $\Gamma$ and $I-\Gamma$ are Gram matrices of linearly independent functions
  on a half circle, so all $N$ parameters lie in $(0,1)$. Hence $V_N$, the
  sum of $p(1-p)$ over the eigenvalues $p$ of $\Gamma$, is
  $\operatorname{tr}\Gamma-\operatorname{tr}\Gamma^2$, and since
  $|\Gamma_{jk}|^2=\pi^{-2}(j-k)^{-2}$ for odd $j-k$ and $\Gamma_{jk}=0$ for
  even $j-k\neq0$,
  \[
    V_N=\frac N4-\frac2{\pi^2}\sum_{\substack{1\leq d<N\\ d\text{ odd}}}
    \frac{N-d}{d^2}
    =\frac2{\pi^2}\sum_{d\text{ odd}}\frac{\min(d,N)}{d^2}
    =\frac{\log N+\gamma+\log2+1}{\pi^2}+O(N^{-2}),
  \]
  using $\sum_{d\text{ odd}}d^{-2}=\pi^2/8$; here $\gamma$ is Euler's
  constant. Rotation by $\pi$ shows that
  $X_N$ and $N-X_N$ have the same law, so $c=N/2$ and $D=N/2+1$. The middle
  form shows that $V_N$ increases with $N$, and a direct evaluation of the
  first form with rational bounds on $\pi$ (a program in the supplement) gives
  $V_{1996}<1<V_{1998}$, so among even $N$ the variance first reaches $1$ at
  $N=1998$. For even $N\geq1998$, Theorem~\ref{thm:main} gives
  $\delta_D\geq1/(4V_N)\sim\pi^2/(4\log N)$, whereas \eqref{eq:ulc} gives
  only $(N+1)/((N/2+2)N/2)\sim4/N$. At $N=10^4$ these lower bounds are about
  $0.215$ and $0.0004$.
\end{example}

The proof of Theorem~\ref{thm:main} runs as follows. The reciprocal bound
\eqref{eq:reciprocal} shows that a small value of $\delta=\delta_D$ forces
small values nearby:
$\delta_{D\pm r}\leq\delta/(1-r\delta)$ while $(r+1)\delta<1$. The masses at distance up to about
$1/\delta$ from the mode are therefore bounded below by explicit multiples of
the maximal mass $M$ (Section~\ref{sec:propagation}). Summed in the pairwise
variance identity, these bounds give $V\geq M^2A(\delta)$ for an explicit
function $A$, while a maximal-mass bound of Bobkov, Marsiglietti, and
Melbourne gives $M^2\geq(1+12V)^{-1}$. Together they give
$V(1+12V)\geq A(\delta)$. Since $V(1+12V)$ increases with $V$ and equals
$(3+\delta)/(4\delta^2)$ at $V=1/(4\delta)$, the theorem reduces to the
one-variable inequality $A(\delta)\geq(3+\delta)/(4\delta^2)$
(Section~\ref{sec:variance}).

Heuristically, the lower bounds on $f_{c\pm r}/M$ are about
$e^{-r^2\delta/2}$ when $r\delta$ is small, so the masses within distance
about $\delta^{-1/2}$ of the mode are comparable to $M$, and $A(\delta)$ is
of order $\delta^{-2}$. With $V(1+12V)\geq A(\delta)$, this makes $V$ of
order at least $1/\delta$.

Section~\ref{sec:scalar} proves the one-variable inequality by Bernstein
expansions. For $1/17\leq\delta<1/4$ the proof rests on 275 rational
coefficients computed by a program and recomputed by an independent checker
(supplement); for $1/22\leq\delta\leq1/17$ on five rational coefficients
printed in the text; and for $\delta\leq1/22$ on coefficients that are
polynomials in an integer parameter, all of whose coefficients are positive.
Several steps can lose a constant, among them the mass bounds, the
maximal-mass bound, the symmetrization in Section~\ref{sec:scalar} that
replaces each $R_r$ by $L_r$, and the replacement of $b_r$ by $\lambda_r$ for
$H\geq16$; they need not be tight simultaneously.

\section{The recurrence for normalized Tur\'an deficits}
\label{sec:propagation}

We show that a small deficit at $D$ forces small deficits nearby and deduce
lower bounds on the masses near the mode.

\subsection{Standing reductions}\label{sec:reduction}

Let $W$ be a finite Bernoulli sum with arbitrary success probabilities
$p_i\in[0,1]$. Put
$I=\{i:0<p_i<1\}$,
$n_1=\#\{i:p_i=1\}$, and
$W^\circ=\sum_{i\in I}B_i$. Then $W=n_1+W^\circ$ and
$\Var(W^\circ)=\Var(W)$. Let $h=(h_k)$ be the pmf of $W^\circ$, where
$h_k=\Pp(W^\circ=k)$. Then $\Pp(W=n_1+k)=h_k$; hence translation shifts the
first-descent index by $n_1$ and
preserves the corresponding adjacent mass ratios and normalized Tur\'an
deficits wherever defined. If $\Var(W)\geq1$, then $I$ is nonempty, so we may
replace $W$ by $W^\circ$, relabel $|I|$ as $n$, and assume $0<p_i<1$ for every
$i$, as in Section~\ref{sec:introduction}.

Under this assumption and $V\geq1$, the set in \eqref{eq:first-descent} is
nonempty. Indeed,
\[
  \frac{f_n}{f_{n-1}}
  =\left(\sum_{i=1}^n\frac{1-p_i}{p_i}\right)^{-1}.
\]
If $f_n\geq f_{n-1}$, then the sum on the right is at most one, whereas
\[
  V=\sum_{i=1}^n p_i^2\frac{1-p_i}{p_i}
  <\sum_{i=1}^n\frac{1-p_i}{p_i}\leq1,
\]
a contradiction. If $f_1<f_0$, then by log-concavity $0$ is a mode, so
$\mathbb EW<1$ by Darroch's theorem \cite[Theorem~4]{darroch1964} and
$V\leq\sum_ip_i=\mathbb EW<1$. Hence $D\geq2$.

If $W'=\sum_{j=1}^mB'_j$ is a Bernoulli sum independent of $W$, then
$W-W'+m=W+\sum_{j}(1-B'_j)$ is a Bernoulli sum with the variance of $W-W'$,
and its pmf is that of $W-W'$ translated by $m$. Hence
Theorem~\ref{thm:main} and Corollary~\ref{cor:support} apply to $W-W'$, with
$D$ the least $k$ in the support of $W-W'$ such that
$\Pp(W-W'=k)<\Pp(W-W'=k-1)$.

\subsection{The recurrence}

For $1\leq k\leq n$, put
\[
  q_k=\frac{f_k}{f_{k-1}}.
\]
Log-concavity says that $(q_k)$ is nonincreasing, and for
$1\leq k\leq n-1$,
\begin{equation}\label{eq:delta-ratio}
  \delta_k=1-\frac{q_{k+1}}{q_k}.
\end{equation}

For the pmf of any finite sum of independent Bernoulli variables, Hillion and
Johnson's Theorem~A.2 and Corollary~A.3 state, for every $k\in\mathbb Z$,
the cubic inequalities \cite[equations~(78)--(79)]{hillionjohnson2016}
\begin{align}
  (f_{k-1}^2-f_{k-2}f_k)f_{k+1}
  &\leq f_{k-1}(f_k^2-f_{k-1}f_{k+1}),
  \label{eq:HJ-left}\\
  (f_{k+1}^2-f_kf_{k+2})f_{k-1}
  &\leq f_{k+1}(f_k^2-f_{k-1}f_{k+1}).
  \label{eq:HJ-right}
\end{align}
Hillion and Johnson extend the pmf by zero outside its support. Since
\[
\begin{aligned}
  f_{k-1}^2-f_{k-2}f_k&=f_{k-1}^2\delta_{k-1},\\
  f_k^2-f_{k-1}f_{k+1}&=f_k^2\delta_k,\\
  f_{k+1}^2-f_kf_{k+2}&=f_{k+1}^2\delta_{k+1},
\end{aligned}
\]
dividing \eqref{eq:HJ-left} by $f_{k-1}f_k^2$ and \eqref{eq:HJ-right} by
$f_{k+1}f_k^2$ gives, for $1\leq k\leq n-1$, the recurrence
\begin{equation}\label{eq:deficit-recurrence}
  \delta_{k-1}(1-\delta_k)\leq\delta_k,
  \qquad
  \delta_{k+1}(1-\delta_k)\leq\delta_k.
\end{equation}
Every $\delta_k$ is positive by \eqref{eq:ulc}, so dividing the second
inequality by $\delta_k\delta_{k+1}$ gives $1/\delta_{k+1}\geq1/\delta_k-1$,
and the first gives the same with $k-1$ in place of $k+1$. The two
inequalities not given by \eqref{eq:deficit-recurrence},
$1/\delta_1-1/\delta_0\geq-1$ and $1/\delta_{n-1}-1/\delta_n\geq-1$, hold
because $1/\delta_0=1/\delta_n=1\leq1/\delta_k$ for every $k$. Hence, for
$0\leq j\leq n-1$,
\begin{equation}\label{eq:reciprocal}
  \left|\frac1{\delta_{j+1}}-\frac1{\delta_j}\right|\leq1.
\end{equation}

Set $\delta=\delta_D$. If $\delta\geq1/4$,
Theorem~\ref{thm:main} follows from $V\geq1$, so it remains to consider
\begin{equation}\label{eq:small-delta}
  0<\delta<\frac14.
\end{equation}
Define
\begin{equation}\label{eq:K-a}
  K=\max\{r\in\mathbb Z_{\geq1}:(r+1)\delta<1\},
  \qquad
  a=\frac{1-2\delta}{1-\delta}.
\end{equation}
By \eqref{eq:small-delta}, $K\geq3$.

\begin{lemma}\label{lem:propagation}
  Assume \eqref{eq:small-delta}. For $0\leq r\leq K$ and each sign,
  $D\pm r\in\{1,\ldots,n-1\}$ and
  \begin{equation}\label{eq:deficit-bound}
    \delta_{D\pm r}\leq\frac{\delta}{1-r\delta}<1.
  \end{equation}
\end{lemma}

\begin{proof}
  We induct on $r$. For $r=0$, we have $D\in\{1,\ldots,n-1\}$ because
  $2\leq D\leq n$ and $\delta_n=1>\delta$, and \eqref{eq:deficit-bound} holds
  with equality. Let $1\leq r\leq K$ and suppose
  $D\pm(r-1)\in\{1,\ldots,n-1\}$. Then $D\pm r\in\{0,\ldots,n\}$, and $r$
  applications of \eqref{eq:reciprocal} give $1/\delta_{D\pm r}\geq1/\delta-r>1$,
  because $(r+1)\delta<1$. So $D\pm r$ is neither $0$ nor $n$, where the
  deficit equals one, and inverting gives \eqref{eq:deficit-bound}.
\end{proof}

By Lemma~\ref{lem:propagation}, $D-K\geq1$ and $D+K\leq n-1$, so every index
$c\pm r$ with $0\leq r\leq K$ lies in the support. Write $M=f_c=\max_k f_k$
for the maximal mass. Since $D\geq K+1\geq2$, the ratio $q_{D-1}$ is
defined, and $q_{D-1}\geq1$ by the definition of $D$.
The lemma with $r=1$ gives $\delta_{D-1}\leq\delta/(1-\delta)$, so
\[
  q_D=q_{D-1}(1-\delta_{D-1})\geq a.
\]

For $0\leq j\leq K-1$, \eqref{eq:deficit-bound} gives
\[
\begin{aligned}
  q_{D+j}
  &=q_D\prod_{s=0}^{j-1}(1-\delta_{D+s})\\
  &\geq a\prod_{s=0}^{j-1}
    \frac{1-(s+1)\delta}{1-s\delta}
   =a(1-j\delta).
\end{aligned}
\]
For $1\leq r\leq K$, we obtain the lower bound
\begin{equation}\label{eq:right-mass-bound}
  R_r:=a^r\prod_{j=1}^{r-1}(1-j\delta),
  \qquad
  \frac{f_{c+r}}M=\prod_{j=0}^{r-1}q_{D+j}\geq R_r.
\end{equation}
On the left, let $1\leq j\leq K$. Then
\[
  q_D=q_{D-j}\prod_{s=1}^j(1-\delta_{D-s})<1,
\]
and hence
\[
\begin{aligned}
  \frac1{q_{D-j}}
  &>\prod_{s=1}^j(1-\delta_{D-s})\\
  &\geq\prod_{s=1}^j
    \frac{1-(s+1)\delta}{1-s\delta}
   =\frac{1-(j+1)\delta}{1-\delta}.
\end{aligned}
\]
For $1\leq r\leq K$, we obtain the lower bound
\begin{equation}\label{eq:left-mass-bound}
  L_r:=(1-\delta)^{-r}
  \prod_{j=2}^{r+1}(1-j\delta),
  \qquad
  \frac{f_{c-r}}M=\prod_{j=1}^r\frac1{q_{D-j}}\geq L_r.
\end{equation}

\section{Variance bounds and scalar reduction}
\label{sec:variance}

We turn these mass bounds into a lower bound on $V$ and reduce
Theorem~\ref{thm:main} to a one-variable inequality. Define auxiliary weights by $w_0=1$, $w_r=R_r$, and $w_{-r}=L_r$ for
$1\leq r\leq K$, and put
\begin{equation}\label{eq:A-def}
  A(\delta)=\frac12\sum_{i,j=-K}^K w_iw_j(i-j)^2.
\end{equation}
Restricting the pairwise variance identity to indices $c-K,\ldots,c+K$, which
lie in the support, and using $f_{c+i}\geq Mw_i$ gives
\begin{equation}\label{eq:pairwise-variance}
  V=\frac12\sum_{i,j=0}^n f_if_j(i-j)^2
  \geq M^2A(\delta).
\end{equation}

We also use the maximal-mass bound of Bobkov, Marsiglietti, and Melbourne
for integer-valued laws
\cite[Corollary~3.2]{bobkovmarsigliettimelbourne2022}:
\begin{equation}\label{eq:maximal-mass-bound}
  V\geq\frac{M^{-2}-1}{12}.
\end{equation}
For completeness, here is a short proof. Let $U$ be independent of $W$ and
uniform on $[-1/2,1/2]$. Then $W+U$ has a density $\varphi$ bounded by $M$,
and $\Var(W+U)=V+1/12$. Put $x_0=\mathbb E(W+U)$ and
$\psi=M\mathbf 1\{|x-x_0|\leq1/(2M)\}$. Both $\varphi$ and $\psi$ integrate
to one, and $(\varphi-\psi)(x)\bigl((x-x_0)^2-(2M)^{-2}\bigr)\geq0$ for
every $x$. Hence
\[
  V+\frac1{12}
  =\int_{\mathbb R}(x-x_0)^2\varphi(x)\,dx
  \geq\int_{\mathbb R}(x-x_0)^2\psi(x)\,dx
  =\frac1{12M^2}.
\]

It remains to prove one scalar estimate.

\begin{proposition}\label{prop:scalar}
  For $0<\delta<1/4$, the quantity in \eqref{eq:A-def} satisfies
  \begin{equation}\label{eq:scalar-target}
    A(\delta)\geq\frac{3+\delta}{4\delta^2}.
  \end{equation}
\end{proposition}

\begin{proof}[Proof of Theorem~\ref{thm:main} from
  Proposition~\ref{prop:scalar}]
  By the reduction preceding \eqref{eq:small-delta}, we may assume
  $0<\delta<1/4$. Suppose that
  $V<1/(4\delta)$. Inequality \eqref{eq:maximal-mass-bound} then gives
  \[
    M^2>\frac{\delta}{3+\delta}.
  \]
  On the other hand, \eqref{eq:pairwise-variance} and
  \eqref{eq:scalar-target} give
  \[
    M^2<\frac1{4\delta A(\delta)}
    \leq\frac{\delta}{3+\delta},
  \]
  a contradiction.
\end{proof}

\section{The scalar inequality}
\label{sec:scalar}

Set
\begin{equation}\label{eq:H-def}
  H=\frac{1-\delta}{\delta}=\frac1\delta-1>3.
\end{equation}
Then $K=\max\{r\in\mathbb Z_{\geq1}:r<H\}$, and the bounds $L_r$ of
\eqref{eq:left-mass-bound} become
\begin{equation}\label{eq:b-def}
  L_r=b_r:=\prod_{s=1}^r\left(1-\frac{s}{H}\right),
\end{equation}
which are nonnegative because $r\leq K<H$. For $1\leq r\leq K$, the bounds
$R_r$ of \eqref{eq:right-mass-bound} satisfy $R_r\geq b_r$. Indeed, if
$C_r=R_r/L_r$, then $C_1=1$ and, for $1\leq r<K$ (so that $H-r-1>0$),
\begin{equation}\label{eq:C-ratio}
  \frac{C_{r+1}}{C_r}
  =\frac{(H-1)(H+1-r)}{(H+1)(H-r-1)}
  =1+\frac{2r}{(H+1)(H-r-1)}\geq1.
\end{equation}

Since every weight is nonnegative, the form \eqref{eq:A-def} is
nondecreasing in each weight: with $K$ fixed,
\[
  \frac{\partial A}{\partial w_\ell}
  =\sum_{j=-K}^K w_j(\ell-j)^2\geq0.
\]
Replacing every $R_r$ by $b_r$ therefore does not increase $A$; denote the
resulting value by $A_{\mathrm{sym}}$. For the symmetric weights $w_0=1$ and
$w_{\pm r}=b_r$, put
\begin{equation}\label{eq:S-T}
  S=1+2\sum_{r=1}^K b_r,
  \qquad
  T=2\sum_{r=1}^K r^2b_r.
\end{equation}
The pairwise identity gives
\[
  A_{\mathrm{sym}}
  =\left(\sum_{r=-K}^K w_r\right)
   \left(\sum_{r=-K}^K r^2w_r\right)
   -\left(\sum_{r=-K}^K rw_r\right)^2
  =ST,
\]
because the final sum vanishes by symmetry. Hence
$A(\delta)\geq A_{\mathrm{sym}}=ST$. Under \eqref{eq:H-def}, the right-hand
side of \eqref{eq:scalar-target} is
\begin{equation}\label{eq:Q-def}
  Q(H)=\frac{(3H+4)(H+1)}4.
\end{equation}
For $H\geq4$, it remains to prove $ST\geq Q(H)$.
Section~\ref{sec:compact} treats $3<H\leq16$ cell by cell, and
Section~\ref{sec:large} treats $H\geq16$.

\subsection{The range \texorpdfstring{$3<H\leq16$}{3 < H <= 16}}
\label{sec:compact}

On $3<H\leq4$ we have $K=3$, and we keep the unsymmetrized bounds $L_r$ and
$R_r$, $r=1,2,3$, because the symmetrized bound fails there: $ST<Q(H)$ at
$H=7/2$, for example. A SymPy computation (in the supplement) gives
\begin{equation}\label{eq:asymmetric-P}
  A(\delta)-Q(H)=\frac{P(H)}{4H^5(H+1)^3},
\end{equation}
where
\begin{align*}
  P(H)={}&-3H^{10}-16H^9+750H^8-3676H^7+6613H^6
  -5460H^5\\
  &+800H^4+4696H^3-7176H^2+4272H-912.
\end{align*}
We substitute $H=3+t$ and expand $P(3+t)$ in the degree-ten Bernstein basis
on $0\leq t\leq1$.

We call each interval $[m,m+1]$ a cell. For $m=4,\ldots,15$ and $H$ in the
cell $[m,m+1]$, define
\begin{equation}\label{eq:compact-ST}
  S_m=1+2\sum_{r=1}^m b_r,
  \qquad
  T_m=2\sum_{r=1}^m r^2b_r,
\end{equation}
and
\begin{equation}\label{eq:compact-P}
  P_m(H)=4H^{2m}\bigl(S_mT_m-Q(H)\bigr).
\end{equation}
For $m<H\leq m+1$, we have $K=m$. At $H=m$, we have $K=m-1$ and $b_m=0$.
Hence $S_m=S$ and $T_m=T$ on the whole cell. Each $P_m$ is an integer
polynomial of degree $2m+2$, and we expand $P_m(m+t)$ in the
degree-$(2m+2)$ Bernstein basis.

If $G(t)=\sum_{j=0}^d a_jt^j$ has degree $d$, the Bernstein conversion used
here is
\begin{equation}\label{eq:bernstein-conversion}
  \beta_i=\sum_{j=0}^i a_j
  \frac{\binom{i}{j}}{\binom{d}{j}},
  \qquad
  G(t)=\sum_{i=0}^d \beta_i\binom di t^i(1-t)^{d-i}.
\end{equation}
The expansions of $P(3+t)$ and of the twelve polynomials $P_m(m+t)$,
$4\leq m\leq15$, have $11$ and $264$ coefficients respectively, and all
$275$ are positive rationals. Since the Bernstein basis polynomials are
nonnegative on $[0,1]$, this proves \eqref{eq:scalar-target} for
$3<H\leq16$. A SymPy program computes each numerator and its Bernstein
coefficients. An independently written checker, using only exact rational
arithmetic from the Python standard library, recomputes them from the
defining formulas. Both programs and the
exact coefficients are in the supplement.

\subsection{The range \texorpdfstring{$H\geq16$}{H >= 16}}
\label{sec:large}

We cover $H\geq16$ by the intervals between consecutive triangular numbers:
choose an integer $J\geq5$ such that
\begin{equation}\label{eq:triangular-cell}
  \frac{J(J+1)}2\leq H\leq\frac{(J+1)(J+2)}2.
\end{equation}
Either choice of $J$ is valid at a shared endpoint. The first interval
($J=5$) is checked directly, and the rest ($J\geq6$) uniformly in $J$. If
$J\geq5$, then $J<J(J+1)/2\leq H$, so $J\leq K$. For $1\leq r\leq J$, the
inequality $\prod_s(1-x_s)\geq1-\sum_sx_s$ for $x_s\in[0,1]$ gives
\begin{equation}\label{eq:bonferroni}
  b_r\geq\lambda_r:=1-\frac{r(r+1)}{2H}\geq0,
\end{equation}
where $\lambda_r\geq0$ because $r(r+1)/2\leq J(J+1)/2\leq H$; this is why
the cells begin at triangular numbers.
Put
\begin{align}
  \tilde S_J&=1+2\sum_{r=1}^J \lambda_r
  =2J+1-\frac{J(J+1)(J+2)}{3H},
  \label{eq:S0}\\
  \tilde T_J&=2\sum_{r=1}^Jr^2\lambda_r
  =\frac{J(J+1)(2J+1)}3-\frac{\sigma_3+\sigma_4}{H},
  \label{eq:T0}
\end{align}
where
\[
  \sigma_3=\sum_{r=1}^Jr^3=\frac{J^2(J+1)^2}{4},
  \qquad
  \sigma_4=\sum_{r=1}^Jr^4=\frac{J(J+1)(2J+1)(3J^2+3J-1)}{30}.
\]
Since $J\leq K$, all omitted weights are nonnegative and $b_r\geq\lambda_r$
for $r\leq J$. Hence $S\geq\tilde S_J$, $T\geq\tilde T_J$, and
$A(\delta)\geq ST\geq\tilde S_J\tilde T_J$. It remains to prove
$\tilde S_J\tilde T_J\geq Q(H)$.

Let $N_J(H)=H^2\bigl(\tilde S_J\tilde T_J-Q(H)\bigr)$. Expanding
\eqref{eq:S0}, \eqref{eq:T0}, and \eqref{eq:Q-def} gives
\begin{align*}
  N_J(H)={}&-\frac34H^4-\frac74H^3
  +\left(\frac{J(J+1)(2J+1)^2}{3}-1\right)H^2\\
  &-\left((2J+1)(\sigma_3+\sigma_4)
    +\frac{J^2(J+1)^2(J+2)(2J+1)}{9}\right)H\\
  &+\frac{J(J+1)(J+2)}{3}(\sigma_3+\sigma_4).
\end{align*}
For $J=5$, only $16\leq H\leq21$ is needed. With $H=16+5t$ and
$0\leq t\leq1$, the degree-four Bernstein coefficients of $N_5(16+5t)$ in $t$
are
\[
  2360,\qquad 7500,\qquad \frac{25055}{2},\qquad
  \frac{254205}{16},\qquad \frac{31115}{2}.
\]

For $J\geq6$, write
\[
  H=\frac{J(J+1)}2+(J+1)t,
  \qquad J=u+6,
  \qquad 0\leq t\leq1.
\]
The degree-four Bernstein coefficients in $t$ of
$N_J\bigl(J(J+1)/2+(J+1)t\bigr)$ are $\beta_i=\mu_i\pi_i(u)/2880$ for
$i=0,\ldots,4$, where
\[
  \begin{gathered}
    \mu_0=J^2(J+1)^2,\qquad \mu_1=J(J+1)^2,\qquad \mu_2=(J+1)^2,\\
    \mu_3=(J+1)^2(J+2),\qquad \mu_4=(J+1)^2(J+2)^2,
  \end{gathered}
\]
and
\begin{align*}
  \pi_0(u)&=121u^4+2474u^3+17431u^2+46014u+25400,\\
  \pi_1(u)&=121u^5+3442u^4+37967u^3+199405u^2+480475u+384510,\\
  \pi_2(u)&=121u^6+4410u^5+65863u^4+512764u^3
    +2172926u^2+4668316u+3829440,\\
  \pi_3(u)&=121u^5+3684u^4+43735u^3+249961u^2+671859u+644400,\\
  \pi_4(u)&=121u^4+2958u^3+25579u^2+88782u+91440.
\end{align*}
Each $\mu_i$ is positive, and each $\pi_i$ has positive coefficients. Thus
$N_J(H)>0$ and $\tilde S_J\tilde T_J>Q(H)$ for $H\geq16$, which with
Section~\ref{sec:compact} proves Proposition~\ref{prop:scalar}. A short
SymPy program in the supplement checks every expansion in this subsection in
exact arithmetic.

\par\bigskip
\begin{supplement}
\stitle{Code and exact certificate data}
\sdescription{The archive, at \url{ZENODO-DOI-PENDING}, contains a SymPy
program that generates all $275$ Bernstein coefficients for $3<H\leq16$, an
independently written checker that recomputes those coefficients using only
the Python standard library, and a SymPy program that checks the expansions
for $H\geq16$, and a program that evaluates $V_{1996}$ and $V_{1998}$ in
Example~\ref{ex:cue} with rational bounds on $\pi$; a manifest gives the
software versions and the commands to rerun the checks. The proof of
Proposition~\ref{prop:scalar} consists of the arguments of
Section~\ref{sec:scalar} together with the exact computations these programs
perform; the deduction of Theorem~\ref{thm:main} from
Proposition~\ref{prop:scalar} (Sections~\ref{sec:propagation}
and~\ref{sec:variance}) is given in full in the text and uses no
computation. The archive also contains a
Lean formalization, conditional on a one-sided form of the recurrence
\eqref{eq:deficit-recurrence} for an abstract nonnegative sequence, of the
deficit bound in Lemma~\ref{lem:propagation} and the exclusion of deficit one
within its window, the bound $q_D\geq a$, and the deduction of \eqref{eq:adjacent-ratio-drop}
from \eqref{eq:main-bound}. Proposition~\ref{prop:scalar}, its computations,
and Theorem~\ref{thm:main} itself are not formalized.}
\end{supplement}

\medskip
\begin{acks}[AI use and author responsibility]
OpenAI's GPT-5.6 model, accessed through Codex, supplied substantial
proof-search assistance for Theorem~\ref{thm:main}. Two ChatGPT (OpenAI)
sessions, run on 3 October 2026 with the model labelled ``Latest'' at Pro
effort, produced a review of an earlier draft and a referee report. The review
proposed the framing in terms of variance and local log-concavity, Corollary~\ref{cor:support}, Proposition~\ref{prop:mode} with
its proof, Proposition~\ref{prop:ulc}, and Example~\ref{ex:cue}; each was
checked, by exact computation on test laws and by independent model-based
proof review, before it was adopted. The referee report pointed out the bound
on $\delta_k$ that follows from Pitman's equation~(20) and its behaviour in
Table~\ref{tab:balanced}, and its independent reconstruction of the
computations of Section~\ref{sec:scalar} and Example~\ref{ex:cue} agreed
exactly with ours. A referee
report from Elicit prompted clarifications of the supplement description and
of Example~\ref{ex:cue}. Claude Opus 5.5 (Anthropic), through Claude Code, assisted
with revision, exact-verification code, and source checking, and other
sessions with Claude, ChatGPT and Codex (both OpenAI), and Gemini (Google)
assisted with computational exploration, proof critique, formalization
support, literature searching, and drafting. Aristotle (Harmonic), an
automated theorem prover, produced the Lean proofs in the conditional
formalization described in the supplement. The author remains responsible for
the mathematics, sources, wording, and released files.
\end{acks}

\bibliographystyle{amsplain}
\bibliography{references}

\end{document}
